\begin{summarybox}
\textbf{30 秒回忆主线：}
\[
\boxed{\text{长轴两端点张角（钝角）}}
\ \longrightarrow\ 
\boxed{\alpha=90^\circ+\arctan\frac{c^2|y|}{2ab^2}}
\ \longrightarrow\ 
\boxed{|y|=b\ \text{时最大}}.
\]

\textbf{必记 4 条：}
\begin{enumerate}
\item \textbf{条件：} $A_1=(-a,0),A_2=(a,0)$，$a>b>0$，
$P\in C\setminus\{A_1,A_2\}$，$c^2=a^2-b^2$。
\item \textbf{最大角：}
$\displaystyle\boxed{\alpha_{\max}=2\arctan(a/b)>90^\circ}$，
仅在 $P=(0,\pm b)$ 取得；
$\displaystyle\boxed{\tan\alpha_{\max}=-2ab/c^2}$。
\item \textbf{全程公式：}
$\displaystyle\boxed{\cot\alpha=-c^2|y|/(2ab^2)}$，
且 $\displaystyle\boxed{k_1k_2=-b^2/a^2}$。
\item \textbf{最小角：}\textbf{没有最小值}，只有
$\displaystyle\boxed{\inf\alpha=90^\circ}$，接近长轴端点时趋近。
\end{enumerate}

\textbf{一眼看穿的等价：}
\[
\boxed{\alpha\ \text{相同}}
\iff\boxed{S=a|y|\ \text{相同}}
\iff\boxed{|y|\ \text{相同}};
\quad
|x|=a\sqrt{1-y^2/b^2}.
\]
通常四个对称点；最大角只有两个短轴端点。

\textbf{最简单的最大值证明：}
先用\textbf{点积/行列式}或\textbf{余弦定理/面积}
推出 $\tan(\alpha-90^\circ)=c^2|y|/(2ab^2)$，
再用 $|y|\le b$；在顶点作角平分线后，
从直角三角形读出 $\tan(\alpha_{\max}/2)=a/b$。

\textbf{换角公式（证明 $\cot\alpha\to\tan(\alpha-90^\circ)$）：}
\[
\boxed{\tan(\alpha-90^\circ)=-\cot\alpha}
\qquad (90^\circ<\alpha<180^\circ).
\]
\textbf{法 A：}用 $\tan=\sin/\cos$，正弦、余弦差角直接得
$-\cos\alpha/\sin\alpha$。
\textbf{法 B：}设 $t=\tan\beta\to+\infty$
（$\beta\to90^{\circ-}$），则
\[
\lim_{t\to+\infty}\frac{\tan\alpha-t}{1+t\tan\alpha}
=\lim_{t\to+\infty}\frac{\tan\alpha/t-1}{1/t+\tan\alpha}
=-\frac1{\tan\alpha}.
\]
\textbf{提醒：}$\tan90^\circ$ 无定义，只有\textbf{趋近}而不能直接代入。

\textbf{C003 对照（切忌混淆）：}
焦点三角形用 $F_1F_2=2c$、面积 $c|y|$、
$\tan(\theta/2)=c|y|/b^2$；
顶点三角形用 $A_1A_2=2a$、面积 $a|y|$、
$\tan(\alpha-90^\circ)=c^2|y|/(2ab^2)$。

\medskip
\textbf{自测 3 问（先口答再看答案）：}
\begin{enumerate}
\item $a=3,b=2$，最大顶角是多少？
\quad \textbf{答：}$\displaystyle2\arctan\frac32$，在 $(0,\pm2)$。
\item 非圆椭圆上存在 $\alpha=80^\circ$ 的点吗？
\quad \textbf{答：}不存在，顶角恒大于 $90^\circ$。
\item 已知某顶角 $\alpha$ 在范围内，$P$ 能唯一确定吗？
\quad \textbf{答：}不能。通常四点，最大角时两点。
\end{enumerate}
\end{summarybox}
